\section{Methods}
This section contains a condensed overview of the neural operator architectures, PDEs, their numerical solution and adjoint calculations, experimental setup, Jacobian sampling, and inversion. Suppmentary materials expand on the details of each component.
\subsection{The SC-NO scheme:} 
Neural Operators (NOs) approximate solution operators for PDEs by learning mappings between function spaces \citep{lu2021learning, kovachki2023neural}. Given an input function \( a(\mathbf{x}, t, p) \)—for example, a spatial field such as an initial condition—NOs learn a parameterized operator \( \mathcal{G}_\theta \) such that \( u(\mathbf{x}, t, p) = \mathcal{G}_\theta(a)(\mathbf{x}, t, p) \), where \( \mathbf{x} \) denotes spatial coordinates, \( t \) denotes time, and \( p \) represents PDE parameters (e.g., viscosity, diffusion coefficients, or forcing magnitudes). The operator thus maps an input function over space--time--parameter domains to an output solution function; for instance, given bed elevation \( a(\mathbf{x}) \), the model predicts water height \( u(\mathbf{x}, t) \) across space and time. SC-NO extends standard NOs by incorporating sensitivity constraints, ensuring accurate gradient predictions alongside state values. The latent representation \( v_t(\mathbf{x}, t, p) \) evolves iteratively as
\( v_{t+1} = \sigma\!\left( \mathcal{K}_\theta(v_t) + \mathcal{L}_\theta(v_t) \right) \),
where \( \mathcal{K}_\theta \) is a nonlocal integral operator and \( \mathcal{L}_\theta \) is a local transformation. The sensitivity-regularized loss minimizes erdo you rors in both state predictions and their parameter sensitivities,
\( \mathcal{L} = \| u - u_{\theta} \|^2 + \lambda \left\| \nabla_{p} u - \nabla_{p} u_{\theta} \right\|^2 \),
where \( \nabla_{p} u \) denotes the gradient of the true solution with respect to the PDE parameters \( p \). Training data—including both solution fields and their sensitivities—are generated using a differentiable PDE solver, enabling consistent supervision of both \( u \) and \( \nabla_{p} u \). SC-NO ensures physics-consistent learning for parametric PDEs, enhancing generalization to unseen conditions. Figure~\ref{fig:sc_no_frame} illustrates the SC-NO framework. A detailed methodological description, along with the rigorous mathematical proof of these guarantees, is provided in Appendix~\ref{sec_sup_method}.


% SC-NOs structure
\begin{figure}[htbp]
    \centering
    \includegraphics[width=0.8\textwidth, trim={0 0 0 0}, clip]{Figuers/SC-NO_framwork.pdf}
    \caption{\scriptsize\textbf{Overview of the Sensitivity-Constrained Neural Operator Framework.} The process begins with input data, including initial states \( a(x, t, p) \) . A neural operator maps these inputs to predicted states \( u(x, t, p) \), while automatic differentiation computes the Jacobian \( \partial u / \partial p \). The framework optimizes a loss function that combines prediction loss \( L_u \) and sensitivity loss \( L_S \), ensuring accuracy and consistency with respect to parameters.}
    \label{fig:sc_no_frame}
\end{figure}



The proposed Sensitivity-Constrained Neural Operator framework is summarized in Figure~\ref{fig:sc_no_frame}. Building on neural operators' universal approximation capabilities, SC-NO incorporates stochastic Jacobian supervision into the training objective by combining the standard prediction loss with a sensitivity regularization term. In each minibatch, the model uses a randomly selected subset of entries from the full Jacobian rather than the entire tensor, and these subsets vary across training so that the model ultimately receives supervision over the full Jacobian without incurring its per-step computational cost. This design improves robustness and generalization to unseen conditions, enabling the model to learn the forward solution operator while simultaneously aligning with parameter sensitivities.

\subsection{Sensitivities through NOs:} 
Either a differentiable numerical model or its neural operator surrogate can compute sensitivities, namely the gradient (or Jacobian) of the output solution with respect to input parameters, forcings, and/or initial conditions. For neural operators, the simplest Jacobian computation is through automatic differentiation (AD), in which we write the forward mode on an AD-supporting platform like PyTorch, JAX, Julia, Tensorflow, etc., and the backward is automatically provided by the platform. We leverage the \texttt{torch.autograd.grad} functionality from PyTorch to efficiently compute Jacobians over different components of the output states, although the same can be provided using other platforms. For acceleration, a proprietary library can parallelize this computation over a minibatch of simulations and multiple output state values, with some cost in increased graphics processing unit (GPU) memory usage. However, AD can rapidly exhaust GPU memory for complex numerical solvers that require the solution of large linear systems or iterative procedures to resolve nonlinearities. In such cases, adjoint-based sensitivity analysis provides a more memory-efficient alternative, with computational cost comparable to a forward--adjoint solve, as described next. While finite-difference methods can be effective in low-dimensional settings, they are not applicable to the high-dimensional cases examined here.


\subsection{Numerical solutions to PDEs and Sensitivity calculations:}
All three PDEs are solved using custom-built, programmatically-differentiable numerical solvers to enable consistent forward simulation and sensitivity evaluation. For PDE1 (ADE) and PDE2 (RANS), the spatially discretized systems are cast as an ordinary differential equation in time and integrated using adaptive Runge-Kutta time stepping, leveraging parts of the \texttt{torchdiffeq} package. Sensitivities for PDE1 are obtained through automatic differentiation of the time integration (not the ``Optimize-then-Discretize" Adjoint offered by \texttt{torchdiffeq}), yielding exact gradients of the discretized dynamics with respect to spatially distributed inputs. For PDE2, sensitivities are computed using a discrete-time adjoint (\textit{Discretize-then-Optimize}) derived from a Lagrangian formulation of the time-discretized RANS system with Spalart–Allmaras closure, enabling efficient gradient evaluation with respect to initial conditions and forcing fields independent of parameter dimensionality.

For PDE3 (Shallow Water), we implemented a finite-volume solver on an unstructured triangular mesh and advance the coupled state variables using an adaptive second-order Heun scheme subject to a CFL stability condition. Numerical fluxes across cell interfaces are computed using the HLLC (Harten-Lax-van Leer-Contact) approximate Riemann solver, similar to the finite volume methods used before \citep{vogl2017high}. Just as for RANS, the corresponding adjoint system is constructed via a discrete-time Lagrangian formulation and solved by backward recursion that mirrors the forward time stepping. Local Jacobians of numerical fluxes and source terms are assembled at the cell level, and gradients with respect to bed topography are obtained by accumulating adjoint-weighted source sensitivities over time. This adjoint formulation computes full-field sensitivities at a cost comparable to a single forward–backward solve, enabling scalable inversion and forecasting for high-dimensional tsunami applications. The Tohoku tsunami case runs a 3,600-second simulation on an unstructured triangular mesh consisting of 308,001 cells and 154,598 nodes, ensuring mass conservation and allowing adaptive local refinement in coastal and impact regions. We use ETOPO 2022 at 30 Arc-Second resolution as initial bathymetry to cover the 1000 km × 1000 km simulation domain (Supplementary Figure \ref{fig:bed_elevation}). The solver has been validated by observed data, e.g., as shown in Supplementary Figure \ref{fig:tohoku_validation} for North Miyagi offshore GPS wave gauge (station 803, NOWPHAS network) in the 2011 Tohoku tsunami event. We have validated the solver on various other small and large events but data are not shown here as the solver is not the focus of this study. The solution and Jacobians are coarsened before training the neural operators.

\subsection{Jacobian sampling and its amortization in the training loop.}
Using reverse-mode differentiation, concurrent AI training can calculate the gradient $\partial u / \partial p$ where $p$ can reach billions or even trillions of parameters, but obtaining the Jacobian entries for many output points $u$ can still be expensive in terms of compute and GPU memory. We therefore sample only a subset of $u$ points in space for the last time point when forming the sensitivity loss. Crucially, the precomputed Jacobian is randomly resampled in every minibatch of the main surrogate training loop. In expectation, the sampled sensitivity loss provides an unbiased estimate of the full Jacobian loss, i.e., $\mathbb{E}[\mathcal{L}{\text{Jac}}^{\text{sampled}}] = \mathcal{L}{\text{Jac}}$. As training proceeds, different Jacobian directions are observed, so that the model effectively covers the full Jacobian space over time while maintaining a bounded per-iteration cost. In this way, Jacobian supervision cost for the SC-FNO part is amortized across training iterations: although only a fraction of sensitivities is enforced at each step, their cumulative effect provides comprehensive gradient-level supervision. We will test sampling more time points in future cases, but in our cases here, sampling the last time points appears sufficient. Unlike prior gradient-informed surrogates that rely on low-dimensional or compressed parameterizations, this approach produces the entire exact Jacobian over training batches.

% \subsection{Okada deformation inversion:} The workflow of the real-world tsunami inversion and forecasting case is composed of three stages. In Stage~1, the initial seafloor deformation field is estimated by gradient-based optimization through the SC-FNO surrogate using the early portion of the buoy time-series data—specifically, the first 30 minutes of a 1-hour event. We compute the loss between synthetic observations and the solutions from the pretrained neural operators (which take the initial deformation and initial water level as inputs and output spatiotemporal water-level dynamics) at the buoy locations. Because the entire workflow supports differentiable programming, we can backpropagate to calculate the gradient of the loss with respect to the seafloor deformation. This deformation field is then adjusted using gradient descent to minimize the loss. Due to the limited information contained in the short span of buoy data, this stage can only provide a rough initial field for the next stage. We use the Okada computational seafloor deformation model to generate data and pretrain a neural network as an encoder mapping the deformation field into nine Okada fault parameters, along with a decoder that maps the parameters back into the deformation field. In Stage~2, we refine the deformation field from Stage~1 using the encoder to produce an initial guess for the Okada parameters used in Stage~3. In Stage~3, the decoder maps the Okada parameters back to the deformation field, which is then fed into the main pretrained neural operators to predict water-level dynamics, and these Okada parameters are optimized again using gradient descent to best match the observed water level. Stages~2 and~3 serve to restrict the inversion to an admissible Okada representation. Because there are only nine parameters in this space, it (i) improves the well-posedness of the inverse problem, given that buoy data are practically limited to sparse locations and a short time window, and (ii) enables the training of accurate encoders and decoders even without gradient supervision. If the full field were observable at all times, the workflow in Stage~1 alone could, in principle, uniquely identify the deformation field. The inferred deformation field is then passed through the pretrained neural operators to predict the full tsunami evolution. Expanded explanations of each stage and the optimization flow are provided in Appendix~\ref{sec:SWE-Tohoku}, Section~\textit{Three-Stage Inversion Framework for Tsunami Source Reconstruction}.
